\section{Стандартная библиотека: итераторы}

\subsection{С какой целью функция, получающая параметр по ссылке, может возвращать тот же параметр тоже по ссылке (pass-through)?}
\begin{definition}[Pass-through]
\textit{Pass-through} — это приём, при котором функция получает объект по ссылке и возвращает ссылку на тот же объект, не создавая копию.
\end{definition}

Такой подход нужен, чтобы:
\begin{itemize}
    \item не копировать объект;
    \item сохранить возможность цепочек вызовов;
    \item передать объект «сквозь» функцию без потери семантики ссылки.
\end{itemize}

\begin{example}
\begin{verbatim}
int& pass(int& x) {
    return x;
}
\end{verbatim}
\end{example}

\begin{remark}
Pass-through полезен, когда функция выступает как промежуточное звено в цепочке преобразований, но не хочет забирать владение объектом и не должна его копировать.
\end{remark}

\subsection{Как условно объекты подразделяют на «большие» и «небольшие»? Какой критерий для этого применяется?}
\begin{definition}[Большие и небольшие объекты]
Объекты условно называют \textbf{небольшими}, если их копирование дешево, и \textbf{большими}, если копирование потенциально дорого.
\end{definition}

Практический критерий — \textbf{правило трёх указателей}: объект считается «большим», если он по смыслу содержит более чем несколько машинных слов, владеет динамической памятью или его копирование требует заметной работы.

\begin{itemize}
    \item \textbf{Небольшие}: \texttt{int}, \texttt{char}, \texttt{double}, маленькие \texttt{enum}.
    \item \textbf{Большие}: \texttt{std::string}, \texttt{std::vector}, крупные структуры, объекты с ресурсами.
\end{itemize}

\subsection{Какие варианты последовательного перебора элементов линейного контейнера существуют?}
Для линейного контейнера, например \texttt{std::vector}, обычно используют три способа:
\begin{itemize}
    \item по индексу через \texttt{operator[]};
    \item циклом \texttt{range-based for};
    \item через итераторы.
\end{itemize}

\begin{example}
\begin{verbatim}
std::vector<int> v{1, 2, 3};

for (std::size_t i = 0; i < v.size(); ++i) {
    std::cout << v[i] << ' ';
}

for (auto x : v) {
    std::cout << x << ' ';
}

for (auto it = v.begin(); it != v.end(); ++it) {
    std::cout << *it << ' ';
}
\end{verbatim}
\end{example}

\subsection{Для чего необходим механизм косвенной адресации, основанный на итераторах?}
\begin{definition}[Итератор]
Итератор — это объект косвенной адресации, который указывает на элемент коллекции и позволяет работать с ним без знания внутреннего устройства контейнера.
\end{definition}

Итераторы нужны для:
\begin{itemize}
    \item унификации доступа к разным контейнерам;
    \item отделения алгоритма от конкретного контейнера;
    \item обхода элементов в стандартных алгоритмах;
    \item абстракции «позиции» в коллекции.
\end{itemize}

\begin{remark}
Итератор связывает контейнер и алгоритм: контейнер хранит элементы, итератор указывает на них, алгоритм их обрабатывает.
\end{remark}

\subsection{Сколько итераторов необходимо для представления одного элемента, двух элементов, n элементов?}
Один итератор указывает на один элемент.  
Но для представления \textbf{диапазона} обычно используют \textbf{два} итератора: начало и конец.

\begin{itemize}
    \item один элемент можно представить как полуинтервал \texttt{[it, it + 1)};
    \item два элемента — \texttt{[it, it + 2)};
    \item \texttt{n} элементов — \texttt{[first, last)}, где расстояние между ними равно \texttt{n}.
\end{itemize}

\subsection{Как с помощью итераторов представляется range (диапазон) объектов?}
\begin{definition}[Range]
\textit{Range} — это совокупность элементов, представляемая парой итераторов в виде полуинтервала \texttt{[first, last)}.
\end{definition}

Здесь:
\begin{itemize}
    \item \texttt{first} указывает на первый элемент диапазона;
    \item \texttt{last} указывает \textbf{на элемент сразу после последнего}.
\end{itemize}

\begin{remark}
Такое представление удобно, потому что пустой диапазон задаётся равенством \texttt{first == last}, а длина диапазона выражается как расстояние между итераторами.
\end{remark}

\subsection{Что общего и что различается у итераторов \texttt{it1} и \texttt{it2} в диапазоне \texttt{[it1, it2)}?}
Общее:
\begin{itemize}
    \item оба относятся к одному диапазону;
    \item оба должны быть совместимы по типу и контейнеру;
    \item между ними определён порядок следования.
\end{itemize}

Различия:
\begin{itemize}
    \item \texttt{it1} — итератор на первый элемент диапазона;
    \item \texttt{it2} — \textit{past-the-end iterator}, то есть указывает \emph{после} последнего элемента;
    \item \texttt{it1} обычно разыменовывается, \texttt{it2} — нет.
\end{itemize}

\subsection{Допустимо ли разыменование \texttt{it1} и \texttt{it2} в диапазоне \texttt{[it1, it2)}?}
\begin{itemize}
    \item \texttt{it1} можно разыменовать, если диапазон не пустой и итератор валиден;
    \item \texttt{it2} разыменовывать нельзя, потому что это итератор «за последним».
\end{itemize}

\begin{remark}
\texttt{it2} — это не элемент коллекции, а специальная граница диапазона. Именно поэтому его называют \textit{sentinel} или \textit{past-the-end iterator}.
\end{remark}

\subsection{Какие операции применимы к подавляющему большинству итераторов? Как умеют «двигаться» итераторы?}
Базовые операции:
\begin{itemize}
    \item конструирование;
    \item копирование;
    \item присваивание;
    \item разыменование \texttt{*it};
    \item сравнение \texttt{it == it2}, \texttt{it != it2};
    \item движение вперёд через \texttt{++it}.
\end{itemize}

Итераторы представляют «позицию» в контейнере и обычно перемещаются на соседний элемент.

\subsection{Чем различаются свободная функция \texttt{std::begin(c)} и метод \texttt{c.begin()}?}
\texttt{c.begin()} — это метод самого контейнера.  
\texttt{std::begin(c)} — свободная функция, которая вызывает подходящий \texttt{begin} для объекта \texttt{c}.

\begin{itemize}
    \item метод удобен для контейнеров с собственным интерфейсом;
    \item свободная функция работает более универсально, в том числе с массивами;
    \item \texttt{std::begin} лучше вписывается в обобщённый код.
\end{itemize}

\begin{example}
\begin{verbatim}
auto it1 = v.begin();
auto it2 = std::begin(v);
\end{verbatim}
\end{example}

\subsection{Какие типы итераторов встроены в стандартные контейнеры? Для чего несколько разновидностей?}
Обычно контейнеры дают:
\begin{itemize}
    \item \texttt{iterator};
    \item \texttt{const\_iterator};
    \item \texttt{reverse\_iterator};
    \item \texttt{const\_reverse\_iterator}.
\end{itemize}

\begin{itemize}
    \item \texttt{iterator} — для чтения и изменения;
    \item \texttt{const\_iterator} — только для чтения;
    \item \texttt{reverse\_iterator} — обход в обратном направлении;
    \item \texttt{const\_reverse\_iterator} — обратный обход без изменения элементов.
\end{itemize}

\begin{remark}
Несколько разновидностей нужны для различной семантики доступа: чтение, запись, прямой обход, обратный обход.
\end{remark}

\subsection{Какие категории итераторов выделяли до стандарта C++20?}
До C++20 обычно выделяли:
\begin{itemize}
    \item \textbf{forward iterators};
    \item \textbf{bidirectional iterators};
    \item \textbf{random access iterators}.
\end{itemize}

\begin{itemize}
    \item forward — движение только вперёд;
    \item bidirectional — вперёд и назад;
    \item random access — произвольный доступ, прыжки на несколько позиций и арифметика итераторов.
\end{itemize}

\begin{remark}
Категории вложены по возрастанию возможностей: random access сильнее bidirectional, а bidirectional сильнее forward.
\end{remark}

\subsection{Какие типы итераторов использует \texttt{std::vector}? А \texttt{std::list}?}
\begin{itemize}
    \item \texttt{std::vector} — random access iterators;
    \item \texttt{std::list} — bidirectional iterators;
    \item многие ассоциативные контейнеры — bidirectional;
    \item некоторые односвязные структуры — forward.
\end{itemize}

\begin{remark}
Чем богаче операции контейнера, тем сильнее категория его итераторов.
\end{remark}

\subsection{Как итерировать от первого к «запоследнему» элементу и обратно? Какие операции при этом используются?}
Прямой обход:
\begin{verbatim}
for (auto it = c.begin(); it != c.end(); ++it) {
    ...
}
\end{verbatim}

Обратный обход:
\begin{verbatim}
for (auto it = c.rbegin(); it != c.rend(); ++it) {
    ...
}
\end{verbatim}

Если есть bidirectional iterator, можно двигаться назад через \texttt{--it}.  
Если есть random access iterator, доступны также \texttt{it += n} и \texttt{it + n}.

\subsection{Движение итераторов: \texttt{++it}, \texttt{it++}, \texttt{it += n}, \texttt{it + n}}
\begin{itemize}
    \item \texttt{++it} — перейти к следующему элементу;
    \item \texttt{it++} — постинкремент, вернуть старое значение и затем сдвинуться;
    \item \texttt{it += n} — сдвиг на \texttt{n} позиций вперёд;
    \item \texttt{it + n} — получить новый итератор, смещённый на \texttt{n}.
\end{itemize}

\begin{remark}
\texttt{++it} обычно предпочтительнее, чем \texttt{it++}, потому что постинкремент может создавать лишнюю копию.
\end{remark}

\subsection{Как сравнивать итераторы при переборе коллекции: \texttt{it != c.end()} или \texttt{it < c.end()}?}
Предпочтительно использовать:
\[
\texttt{it != c.end()}
\]

Потому что:
\begin{itemize}
    \item это поддерживается большинством категорий итераторов;
    \item это естественный способ проверки конца диапазона;
    \item \texttt{<} доступно не для всех итераторов.
\end{itemize}

\begin{remark}
Сравнение через \texttt{<} допустимо только для random access iterator, например у \texttt{std::vector}. Но универсальный код должен опираться на \texttt{!=}.
\end{remark}

\subsection{Можно ли копировать итератор \texttt{c.begin()}? Можно ли его разыменовать? Можно ли двигать вперёд?}
Да, копировать можно.  
Да, если он указывает на валидный элемент, его можно разыменовать.  
Да, можно двигать вперёд, если контейнер и категория итератора это позволяют.

\begin{remark}
Итератор — это полноценный объект, который удобно передавать по значению.
\end{remark}

\subsection{Можно ли копировать итератор \texttt{c.end()}? Можно ли его разыменовать? Можно ли двигать назад?}
Копировать \texttt{c.end()} можно.  
Разыменовывать нельзя.  
Двигать назад можно только если категория итератора это допускает, например для bidirectional или random access.

\begin{remark}
\texttt{end()} — это допустимая граница диапазона, но не элемент диапазона.
\end{remark}

\subsection{Какие операции допустимы к итератору на валидный элемент и к итератору \textit{past-the-end}?}
Для итератора на валидный элемент обычно допустимы:
\begin{itemize}
    \item разыменование;
    \item инкремент;
    \item сравнение;
    \item в зависимости от категории — декремент и прыжки.
\end{itemize}

Для \textit{past-the-end} итератора:
\begin{itemize}
    \item разыменование запрещено;
    \item сравнение допустимо;
    \item движение назад возможно только для соответствующих категорий и только если это не нарушает валидность.
\end{itemize}

\subsection{Какой тип итератора следует выбирать при разработке нового алгоритма: минимально достаточный или максимально сильный?}
Следует требовать \textbf{минимально достаточную} категорию итератора, которая реально нужна алгоритму.

\begin{itemize}
    \item если алгоритму достаточно однопроходного обхода, не надо требовать random access;
    \item если нужен только просмотр вперёд, не надо требовать bidirectional;
    \item чем слабее требование, тем больше контейнеров сможет использовать алгоритм.
\end{itemize}

\begin{remark}
Это важный принцип обобщённого программирования: алгоритм должен быть как можно менее требовательным.
\end{remark}

\subsection{Что происходит с итератором после удаления элемента из коллекции по позиции этим итератором?}
Итератор, по которому удалили элемент, обычно становится невалидным.

\begin{itemize}
    \item после \texttt{erase(it)} старый \texttt{it} использовать нельзя;
    \item обычно \texttt{erase} возвращает новый валидный итератор на следующий элемент;
    \item его нужно сохранить и использовать дальше.
\end{itemize}

\begin{example}
\begin{verbatim}
auto it = v.erase(it); // старый it уничтожен, используется возвращённый
\end{verbatim}
\end{example}

\subsection{Краткие ответы на контрольные вопросы}
\begin{itemize}
    \item \textbf{Зачем pass-through?} Чтобы вернуть ссылку на тот же объект без копирования.
    \item \textbf{Что такое iterator?} Объект косвенной адресации и связующий элемент между контейнером и алгоритмом.
    \item \textbf{Сколько итераторов нужно для диапазона?} Обычно два: начало и конец.
    \item \textbf{Почему \texttt{[first, last)}?} Потому что конец не включается, это упрощает пустые диапазоны и алгоритмы.
    \item \textbf{Что такое \texttt{end()}?} Итератор на позицию после последнего элемента.
    \item \textbf{Почему \texttt{it != end()} лучше, чем \texttt{it < end()}?} Потому что универсальнее.
    \item \textbf{Какие категории итераторов существуют?} Forward, bidirectional, random access.
    \item \textbf{Что делать с итератором после \texttt{erase}?} Считать старый итератор недействительным и использовать возвращённый.
\end{itemize}
